\documentclass[12pt]{article}
\usepackage{amsmath, amssymb}
\usepackage[a4paper, margin=1in]{geometry}

\begin{document}

\begin{center}
    {\LARGE \textbf{Lecture 19: N-Variate Joint PDF / PMF / CDF}}
\end{center}

\section*{1. Joint Distribution of N Random Variables}

\subsection*{Definition}

For random variables $(X_1, X_2, \dots, X_n)$, we define the vector:

\[
\mathbf{X} =
\begin{bmatrix}
X_1 \\
X_2 \\
\vdots \\
X_n
\end{bmatrix}
\]

\begin{itemize}
    \item The \textbf{joint distribution} describes the probability behavior of all variables \textbf{together}.
    \item It can be represented using:
    \begin{itemize}
        \item Joint PMF (for discrete variables)
        \item Joint PDF (for continuous variables)
        \item Joint CDF
    \end{itemize}
\end{itemize}

\section*{2. Joint Gaussian Distribution}

\subsection*{Joint PDF for Two Variables}

For two random variables $(X)$ and $(Y)$:

\[
f_{XY}(x,y) =
\frac{1}{2\pi \sigma_X \sigma_Y \sqrt{1 - \rho^2}}
\exp \left\{
-\frac{1}{2(1-\rho^2)} \left[
\left(\frac{x-\mu_X}{\sigma_X}\right)^2
+
\left(\frac{y-\mu_Y}{\sigma_Y}\right)^2
- 2\rho \left(\frac{x-\mu_X}{\sigma_X}\right)
\left(\frac{y-\mu_Y}{\sigma_Y}\right)
\right]
\right\}
\]

\subsection*{Parameters}

\begin{itemize}
    \item $(\mu_X, \mu_Y)$: Means
    \item $(\sigma_X, \sigma_Y)$: Standard deviations
    \item $(\rho)$: Pearson correlation coefficient
\end{itemize}

\section*{3. Example 1: Jointly Gaussian Variables}

\subsection*{Given}

\begin{itemize}
    \item $(\mu_X = 2, \sigma_X = 1)$
    \item $(\mu_Y = -3, \sigma_Y = 2)$
    \item $\text{cov}(X,Y) = -1$
\end{itemize}

\subsection*{(a) Joint PDF}

\subsubsection*{Step 1: Find Correlation Coefficient}

\[
\rho_{XY} = \frac{\text{cov}(X,Y)}{\sigma_X \sigma_Y}
\]

Substitute values:
\[
\rho_{XY} = \frac{-1}{1 \cdot 2} = -0.5
\]

\subsubsection*{Step 2: Substitute into Joint PDF}

Using:

\begin{itemize}
    \item $(\mu_X = 2)$
    \item $(\mu_Y = -3)$
    \item $(\sigma_X = 1)$
    \item $(\sigma_Y = 2)$
    \item $(\rho = -0.5)$
\end{itemize}

Final joint PDF is obtained by substituting into the general formula.

\subsection*{(b) Marginal PDFs}

\subsubsection*{Definition}

\[
f_X(x) = \int_{-\infty}^{\infty} f_{XY}(x,y)\,dy
\]
\[
f_Y(y) = \int_{-\infty}^{\infty} f_{XY}(x,y)\,dx
\]

\subsubsection*{Key Property (Given in PPT)}

If variables are jointly Gaussian:

\begin{itemize}
    \item $X \sim \mathcal{N}(\mu_X, \sigma_X^2)$
    \item $Y \sim \mathcal{N}(\mu_Y, \sigma_Y^2)$
\end{itemize}

Thus:

\begin{itemize}
    \item $X \sim \mathcal{N}(2,1^2)$
    \item $Y \sim \mathcal{N}(-3,2^2)$
\end{itemize}

\subsection*{(c) Find $\Pr(X < 0)$}

\subsubsection*{Step 1: Use CDF}

\[
F_X(x) = \Pr(X < x)
\]
\[
\Pr(X < 0) = F_X(0)
\]

\subsubsection*{Step 2: Standardization}

Convert to standard normal $(Z \sim \mathcal{N}(0,1))$:

\[
Z = \frac{X - \mu_X}{\sigma_X}
\]

\[
\Pr(X < 0) = \Pr\left(Z < \frac{0 - 2}{1}\right)
= \Pr(Z < -2)
\]

\subsubsection*{Step 3: Symmetry Property}

\[
\Phi(-x) = 1 - \Phi(x) = \Pr(Z > x)
\]

\[
\Pr(Z < -2) = \Pr(Z > 2)
\]

\subsubsection*{Step 4: Q-function}

\[
Q(x) = \Pr(Z > x)
\]

\[
\boxed{\Pr(X < 0) = Q(2)}
\]

\section*{4. Application: Large MIMO System}

\subsection*{System Description}

\begin{itemize}
    \item Wireless communication system with:
    \begin{itemize}
        \item $(M)$ transmit antennas
        \item $(N)$ receive antennas
    \end{itemize}
\end{itemize}

\subsection*{Model}

\begin{itemize}
    \item Each channel path between antennas is modeled as a \textbf{random variable}:
    \[
    h_{ij}
    \]
\end{itemize}

Thus, system behavior involves \textbf{multiple random variables jointly}, requiring joint distributions.

\section*{5. Additional Topics Mentioned}

\subsection*{Joint PMF / PDF / CDF}

\begin{itemize}
    \item Used to represent joint behavior of multiple variables.
\end{itemize}

\subsection*{Independence of Random Variables}

\begin{itemize}
    \item Mentioned as a topic (no derivation shown in provided content).
\end{itemize}

\section*{6. Other Examples (Listed in Outline)}

\begin{itemize}
    \item Example 2: Uniformly distributed 3D vector
    \item Example 3: Trivariate random variables
\end{itemize}

(Note: No detailed derivations present in the provided content.)


\end{document}